\section{短时序建模：单帧、早融合、晚融合和 3D CNN}

\subsection{单帧 baseline 为什么常常不差}

最朴素的做法，是把视频当成一堆图像来处理：抽取单帧、丢给图像分类器、最后在时间上平均或投票。这个 baseline 常常 surprisingly strong，因为很多动作本身就带着明显的静态线索，比如场景、道具、人体姿态和运动装备。

\begin{importantbox}{单帧 baseline 的意义}
\begin{itemize}
    \item 它是最便宜的对照组。
    \item 它能告诉你“外观信息到底有多强”。
    \item 如果单帧已经很强，说明任务可能被静态背景严重污染。
\end{itemize}
\end{importantbox}

\subsection{Early fusion, late fusion, 3D CNN}

短时序建模阶段最重要的不是某一个单独网络，而是“什么时候开始把时间维度纳入计算”。课程把它分成三类：

\begin{center}
\begin{tabular}{p{2.8cm} p{4.3cm} p{4.0cm}}
\toprule
方法 & 核心做法 & 主要特点 \\
\midrule
Late fusion & 每帧先用 2D CNN 编特征，最后再聚合时间 & 先学空间，再合并时间，最朴素也最稳妥 \\
Early fusion & 一开始就把多帧堆成更厚的 channel & 早早混合时间信息，但很粗暴 \\
3D CNN & 用 3D 卷积和 3D pooling 逐步融合时空信息 & 更自然，但算力更贵 \\
\bottomrule
\end{tabular}
\end{center}

\begin{figure}[H]
\centering
\includegraphics[width=0.95\textwidth]{slides-images/slide-020.jpg}
\caption{Early fusion：把时间堆成 channel，在第一层 2D 卷积里一次性看到所有帧\protect\footnotemark}
\end{figure}
\footnotetext{Slides 第21页。}

\begin{figure}[H]
\centering
\includegraphics[width=0.95\textwidth]{slides-images/slide-021.jpg}
\caption{3D CNN：把卷积和 pooling 都扩展到时间维，让网络从一开始就有时空归纳偏置\protect\footnotemark}
\end{figure}
\footnotetext{Slides 第22页。}

\subsection{为什么 3D CNN 是“慢融合”}

3D CNN 的直觉很好理解：如果视频里的模式本来就会在时间上移动，那卷积也应该在时间上滑动。于是，输入从 $C \times H \times W$ 变成 $C \times T \times H \times W$，卷积核也从 $k_H \times k_W$ 变成 $k_T \times k_H \times k_W$。

\begin{knowledgebox}{3D 卷积的基本形式}
如果输入特征是 $X \in \mathbb{R}^{C \times T \times H \times W}$，那么一个 3D 卷积核会在 $(T, H, W)$ 三个维度上滑动，输出仍然保留时间维。这样做的效果是：模型在每一层都能逐渐扩大时空感受野，而不是把时间只留到最后一层处理。
\end{knowledgebox}

\begin{figure}[H]
\centering
\includegraphics[width=0.95\textwidth]{slides-images/slide-022.jpg}
\caption{从 2D 卷积自然过渡到 3D 卷积：时间维和空间维被统一看成一个局部立方体\protect\footnotemark}
\end{figure}
\footnotetext{Slides 第23页。}

\begin{figure}[H]
\centering
\includegraphics[width=0.95\textwidth]{slides-images/slide-024.jpg}
\caption{Late fusion 的 toy example：时间维先保留到最后，卷积主要先在空间上展开\protect\footnotemark}
\end{figure}
\footnotetext{Slides 第25页。}

\begin{figure}[H]
\centering
\includegraphics[width=0.95\textwidth]{slides-images/slide-027.jpg}
\caption{Early fusion 的 toy example：时间信息在最早阶段被压成 channel，空间特征之后再逐步提取\protect\footnotemark}
\end{figure}
\footnotetext{Slides 第28页。}

\begin{figure}[H]
\centering
\includegraphics[width=0.95\textwidth]{slides-images/slide-030.jpg}
\caption{3D CNN 的 toy example：时间和空间一起逐步融合，因此也叫 slow fusion\protect\footnotemark}
\end{figure}
\footnotetext{Slides 第31页。}

\begin{importantbox}{三种融合方式的核心差别}
\begin{itemize}
    \item Late fusion 是“先空间、后时间”。
    \item Early fusion 是“先时间、后空间”。
    \item 3D CNN 是“时间和空间一起慢慢融合”。
\end{itemize}
\end{importantbox}

\subsection{为什么 early fusion 还不够}

讲者在 toy example 里给出的对比很清楚：early fusion 的第一层 2D 卷积虽然看到了多帧，但它本质上还是把时间当作额外 channel，一次性压缩。问题在于，一层卷积通常不足以编码复杂动作的全部时间结构。很多动作不是“某一帧像什么”，而是“变化过程如何发生”。

\begin{warningbox}{early fusion 的局限}
它把时间信息“早早压扁”，因此对需要逐层抽象运动模式的任务不够灵活。这个设计在工程上简单，但表达能力有限。
\end{warningbox}

\subsection{C3D：3D CNN 的 VGG 风格版本}

3D CNN 真正成为成熟方案，C3D 是里程碑之一。它像 VGG 一样堆很多简单的 `3x3x3` 卷积和 `2x2x2` pooling，把视频看成一个真正的 3D 体。其结构简洁，便于迁移，也容易作为 feature extractor 直接下游使用。

\begin{figure}[H]
\centering
\includegraphics[width=0.95\textwidth]{slides-images/slide-039.jpg}
\caption{C3D 是 3D CNN 的经典骨架：像 VGG 一样堆叠标准化的 3D 卷积块\protect\footnotemark}
\end{figure}
\footnotetext{Slides 第40页。}

\begin{figure}[H]
\centering
\includegraphics[width=0.95\textwidth]{slides-images/slide-040.jpg}
\caption{C3D 的层级结构：从 16 帧输入一路下采样到全连接层，形成可迁移的视频特征提取器\protect\footnotemark}
\end{figure}
\footnotetext{Slides 第41页。}

\begin{figure}[H]
\centering
\includegraphics[width=0.95\textwidth]{slides-images/slide-041.jpg}
\caption{C3D 的计算代价不低：3x3x3 的 3D 卷积非常贵，因此效率成为视频模型的重要约束\protect\footnotemark}
\end{figure}
\footnotetext{Slides 第42页。}

\begin{knowledgebox}{C3D 的重要性}
C3D 的价值不在于它是最终答案，而在于它把“视频也可以像图像一样用标准卷积堆起来”这件事讲清楚了。它把视频理解从概念问题推进到可工程化实现的问题。
\end{knowledgebox}

\begin{figure}[H]
\centering
\includegraphics[width=0.95\textwidth]{slides-images/slide-042.jpg}
\caption{Sports-1M 上的对比：3D CNN 和 C3D 都比单帧更强，但代价也明显更高\protect\footnotemark}
\end{figure}
\footnotetext{Slides 第43页。}

\begin{importantbox}{短时序阶段的结论}
\begin{enumerate}
    \item 单帧 baseline 是必须先看的参照。
    \item Early fusion 简单，但时间压缩太早。
    \item 3D CNN 更自然地融合时空信息。
    \item 但 3D 卷积计算代价高，后续方法必须继续追求效率。
\end{enumerate}
\end{importantbox}

